% =============================================================================
%  Capítulo 4 — Resultados
%
%  En tu PDF actual este es el Capítulo 5. Renumerar al borrar el "Cap. 4
%  Desarrollo" vacío.
%
%  Una subsección por requisito del enunciado:
%   §4.1 ⟷ requisito 1 (procesos VE/VP)
%   §4.2 ⟷ requisito 3 (schedules)
%   §4.3 ⟷ requisito 2 (samplers)        ← FALTA en tu PDF
%   §4.4 ⟷ requisito 5 (controllable)    ← FALTA imputación
%   §4.5 ⟷ requisito B&N + color
%   §4.6 ⟷ rendimiento (RNF)
% =============================================================================
\chapter{Resultados}
\label{cap:resultados}

\todo{Introducción breve del capítulo: qué se compara, sobre qué datasets,
con qué métricas, condiciones de evaluación reproducibles
(seed, hardware, número de muestras).}

% -----------------------------------------------------------------------------
\section{Comparación de procesos: VE vs VP}
\label{sec:res-ve-vp}

\todo{Setup: mismo dataset (MNIST), mismo sampler (Euler--Maruyama),
misma seed, mismo número de pasos $N$. Variar solo el proceso.

Incluir:
\begin{itemize}
  \item Galería visual lado a lado (8--16 muestras por proceso).
  \item Tabla con FID, IS y BPD por proceso.
  \item Curvas de loss vs época superpuestas.
  \item Discusión: estabilidad de entrenamiento, calibración de
        $\sigma_\text{max}$ (VE) vs $\beta_\text{max}$ (VP).
\end{itemize}}

% \begin{figure}[H]
%   \centering
%   \includegraphics[width=0.95\textwidth]{galeria_ve_vp.pdf}
%   \caption{Comparativa visual VE (arriba) vs VP (abajo) sobre MNIST.}
%   \label{fig:ve-vs-vp}
% \end{figure}

% \begin{table}[H]
%   \centering
%   \begin{tabular}{lccc}
%     \toprule
%     Proceso & FID $\downarrow$ & IS $\uparrow$ & BPD $\downarrow$ \\
%     \midrule
%     VE & TODO & TODO & TODO \\
%     VP & TODO & TODO & TODO \\
%     \bottomrule
%   \end{tabular}
%   \caption{Métricas VE vs VP sobre MNIST.}
% \end{table}

% -----------------------------------------------------------------------------
\section{Comparación de schedules: lineal vs coseno (sobre VP)}
\label{sec:res-schedules}

\todo{Setup: VP fijo, mismo sampler, mismo dataset.

Incluir:
\begin{itemize}
  \item Galería visual.
  \item Trayectoria forward (figura tipo Nichol \& Dhariwal pero
        \emph{generada con tu código}, no copiada).
  \item Tabla FID/IS/BPD.
  \item Discusión: por qué coseno conserva señal hasta más tarde y
        ayuda en baja resolución.
\end{itemize}}

% -----------------------------------------------------------------------------
\section{Comparación de samplers: EM vs PC vs PF-ODE}
\label{sec:res-samplers}

\todo{NUEVA respecto a tu PDF. Es la sección más vendible del capítulo.

Setup: mismo modelo entrenado, misma seed, distinto sampler.

Incluir:
\begin{itemize}
  \item Galería visual lado a lado.
  \item Tabla con FID, tiempo por imagen y \#evals de red por muestra.
  \item \textbf{Curva calidad vs N} (FID en función de
        $N \in \{10, 25, 50, 100, 250, 500, 1000\}$). Esta gráfica
        es la que mejor muestra la diferencia entre samplers.
  \item Demostración del determinismo del PF-ODE: dos corridas con la
        misma seed dan resultado idéntico; con EM o PC, distintas.
\end{itemize}}

% -----------------------------------------------------------------------------
\section{Generación controlable}
\label{sec:res-controlable}

\subsection{Class-conditional con CFG}
\label{sec:res-cfg}

\todo{Rejilla 10$\times$8 sobre SVHN: filas = clase, columnas = muestras.

\textbf{Sweep del peso de guidance} $w \in \{0, 1, 3, 5, 10\}$ para mostrar
el trade-off fidelidad-clase vs diversidad.}

\subsection{Imputación}
\label{sec:res-imputacion}

\todo{NUEVA respecto a tu PDF. Requisito 5b del enunciado.

Tres filas por ejemplo:
\begin{itemize}
  \item Imagen original.
  \item Máscara aplicada.
  \item Reconstrucción (varias muestras para mostrar diversidad).
\end{itemize}
Probar al menos tres tipos de máscara: bandas horizontales, cuadrado
central, ruido aleatorio.}

% -----------------------------------------------------------------------------
\section{Resultados sobre dataset color}
\label{sec:res-color}

\todo{Demuestra el "B\&W or color images" del enunciado.

Sobre los gatos / SVHN:
\begin{itemize}
  \item Galería de muestras de los checkpoints (5, 50, 300 épocas...).
  \item Métricas FID/IS para cada checkpoint.
  \item Discusión: relación con el fenómeno del double descent
        (si se mantiene esa observación en §2).
\end{itemize}}

% -----------------------------------------------------------------------------
\section{Análisis de rendimiento}
\label{sec:res-rendimiento}

\todo{Verifica los Requisitos No Funcionales cuantitativos del
§\ref{sec:rnf}.

Tabla:
\begin{itemize}
  \item Filas: combinaciones (proceso $\times$ sampler $\times$ $N$).
  \item Columnas: segundos/imagen, memoria pico GPU, \#evals red,
        cumple-RNF (sí/no).
\end{itemize}}
